布朗运动的独立增量与高斯结构可把路径计算化为协方差积分。确定的平方可积核对应高斯伊藤积分；指数化或在独立随机时刻取值后，一般不再正态。联合高斯条件可用正交投影处理，最大值可用反射原理；对终点等零概率条件，须通过密度或连续条件版本解释。

独立指数失效可转成贴现因子，再用生成元边值方程和伊藤公式验证。交换积分、期望或参数导数时，应说明非负性、可积性或支配条件。泊松分流要求独立标记；更新过程在长时间区间内均匀选点观察，其极限分布须按间隔长度加权。
\par\smallskip
